\documentclass[11pt,a4paper]{article}

\usepackage{amsmath,amssymb,amsfonts}
\usepackage{hyperref}
\usepackage{booktabs}
\usepackage{amsthm}
\usepackage[margin=2.5cm]{geometry}

\newtheorem{theorem}{Theorem}
\newtheorem{proposition}[theorem]{Proposition}
\newtheorem{conjecture}[theorem]{Conjecture}
\newtheorem{corollary}[theorem]{Corollary}

\newcommand{\Spec}{\mathrm{Spec}}
\newcommand{\Zp}{\mathbb{Z}[1/p]}
\newcommand{\Qp}{\mathbb{Q}_p}
\newcommand{\Zhat}{\hat{\mathbb{Z}}}

\begin{document}

\title{Cosmological Implications of Arithmetic Spacetime:\\
  If the Universe Lives on $\Spec(\mathbb{Z})$}

\author{Wright Brothers\\{\small Independent Research}}
\date{April 2026}
\maketitle

\begin{abstract}
The arithmetic vacuum engineering program rests on a
foundational premise: the quantum vacuum has the arithmetic
structure of $\Spec(\mathbb{Z})$, the scheme of integers.
We examine the cosmological implications of this premise.
If spacetime is fundamentally arithmetic, then:
(1)~the cosmological constant equals $\zeta(-3)/120$ in Planck
units, potentially resolving the cosmological constant problem
by replacing cutoff regularization with $\zeta$-regularization
as the \emph{correct} description;
(2)~the Big Bang is a phase transition of the Bost-Connes system
at $\beta = 1$, where the Galois symmetry
$\mathrm{Gal}(\mathbb{Q}^{\mathrm{ab}}/\mathbb{Q})$ spontaneously
breaks, differentiating the primes;
(3)~dark energy is the Casimir energy of $\Spec(\mathbb{Z})$ itself,
i.e., the zero-point energy of the arithmetic vacuum;
(4)~the fine structure constant and other dimensionless constants
are determined by the arithmetic geometry of $\Spec(\mathbb{Z})$
(zeta values, class numbers, regulator volumes);
(5)~the multiverse, if it exists, is the set of localizations
$\Spec(\mathbb{Z}[S^{-1}])$ for different sets $S$ of primes,
each defining a universe with different physical laws
via topos logic change.
We formulate each implication as a precise conjecture and
identify the observational signatures that would distinguish
arithmetic spacetime from conventional cosmology.
\end{abstract}

% ============================================================================
\section{Introduction}
\label{sec:intro}
% ============================================================================

The arithmetic vacuum engineering program~\cite{paper_A,paper_B,paper_C}
is built on a premise that is rarely stated explicitly:

\begin{conjecture}[Arithmetic spacetime]
\label{conj:arith}
The fundamental structure of spacetime is the arithmetic scheme
$\Spec(\mathbb{Z})$. The quantum vacuum is a sheaf on this scheme.
The Riemann zeta function $\zeta(s)$, as the zeta function of
$\Spec(\mathbb{Z})$, encodes the vacuum's physical properties.
\end{conjecture}

This paper does not prove the conjecture. Instead, we ask:
\emph{if it is true, what are the consequences for cosmology?}

The results are striking. Several of the deepest puzzles in
modern cosmology --- the cosmological constant problem, the
origin of the Big Bang, the nature of dark energy, the values
of fundamental constants --- acquire natural (if speculative)
explanations within the arithmetic framework.

% ============================================================================
\section{The cosmological constant problem}
\label{sec:cc}
% ============================================================================

\subsection{The problem}

The cosmological constant $\Lambda$ is the energy density of
the vacuum. Quantum field theory predicts
\begin{equation}
  \rho_{\mathrm{vac}}^{\mathrm{QFT}} \sim \frac{\hbar c}{l_P^4}
  \sim 10^{113}\,\mathrm{J/m^3},
\end{equation}
while observation gives
$\rho_{\mathrm{vac}}^{\mathrm{obs}} \sim 10^{-9}\,\mathrm{J/m^3}$.
The discrepancy of $\sim 10^{122}$ is the worst prediction in
the history of physics~\cite{Weinberg1989}.

\subsection{The arithmetic resolution}

The QFT prediction uses \emph{cutoff regularization}:
$\rho \sim \int_0^{E_P} \omega^3\,d\omega \sim E_P^4$.
This diverges and must be cut off at the Planck scale.

Zeta regularization replaces this with
\begin{equation}
  \rho_{\mathrm{vac}}^{\zeta} \propto \zeta(-3) = \frac{1}{120}.
  \label{eq:cc_zeta}
\end{equation}

\begin{conjecture}[Arithmetic cosmological constant]
\label{conj:cc}
The physical cosmological constant is
\begin{equation}
  \Lambda = \frac{8\pi G}{c^4} \cdot \frac{E_P}{l_P^3} \cdot \zeta(-3)
  = \frac{8\pi}{120}\,l_P^{-2},
\end{equation}
where the factor $\zeta(-3) = 1/120$ replaces the divergent
cutoff integral.
\end{conjecture}

This gives $\Lambda \sim 0.2\,l_P^{-2}$, which is still far
from the observed value $\Lambda_{\mathrm{obs}} \sim 10^{-122}\,l_P^{-2}$.
However, the key point is not the numerical value but the
\emph{framework}: if the vacuum is $\Spec(\mathbb{Z})$, then
the cosmological constant is determined by $\zeta(-3)$, not by
a cutoff. The remaining discrepancy may be explained by
the \emph{partial muting} of primes in the actual vacuum
(see Section~\ref{sec:dark}).

% ============================================================================
\section{The Big Bang as Bost-Connes phase transition}
\label{sec:bigbang}
% ============================================================================

The Bost-Connes system has a phase transition at $\beta = 1$:

\begin{itemize}
\item $\beta < 1$ (high temperature): unique KMS state.
  Full symmetry $\mathrm{Gal}(\mathbb{Q}^{\mathrm{ab}}/\mathbb{Q})$
  is preserved. All primes are indistinguishable.
\item $\beta > 1$ (low temperature): the symmetry spontaneously
  breaks. KMS states form a continuous family parameterized by
  $\Zhat^* = \prod_p \mathbb{Z}_p^*$. Primes become individually
  addressable.
\end{itemize}

\begin{conjecture}[Arithmetic Big Bang]
\label{conj:bigbang}
The Big Bang is the Bost-Connes phase transition at $\beta = 1$.

Before the transition ($\beta < 1$): the universe is in the
symmetric phase. No individual primes exist. No particle species
are distinguishable (because the Standard Model requires
the broken phase of the internal space $F$).
Spacetime has maximal symmetry.

At the transition ($\beta = 1$): the Galois symmetry breaks.
Primes differentiate. The Euler product structure of $\zeta(s)$
emerges. Particle species appear (via Connes' spectral action
on $M^4 \times F$). The arrow of time appears (because the
broken phase selects a specific KMS state).

After the transition ($\beta > 1$): the universe we observe.
Each prime $p$ contributes its Euler factor $(1-p^{-\beta})^{-1}$
to the vacuum structure.
\end{conjecture}

\textbf{Observational signature.}
The cosmic microwave background (CMB) encodes fluctuations from
the epoch near the phase transition. If the Big Bang is a BC
transition, the primordial power spectrum should reflect the
number-theoretic structure of $\zeta(s)$ near $\beta = 1$:
\begin{equation}
  P(k) \propto \left|\frac{\zeta(1 + \epsilon + ik)}{\zeta(1+\epsilon)}\right|^2
\end{equation}
for small $\epsilon > 0$ (proximity to the critical point).
The Riemann zeros on the critical line would imprint as
oscillatory features in the CMB power spectrum.

% ============================================================================
\section{Dark energy as arithmetic Casimir energy}
\label{sec:dark}
% ============================================================================

\begin{conjecture}[Arithmetic dark energy]
\label{conj:dark}
Dark energy is the Casimir energy of $\Spec(\mathbb{Z})$ itself:
the zero-point energy of the vacuum computed on the arithmetic
scheme, with boundary conditions set by the large-scale topology
of the universe.
\end{conjecture}

In standard Casimir physics, the vacuum energy depends on
boundary conditions (plate separation, geometry).
On $\Spec(\mathbb{Z})$, the ``boundaries'' are:
\begin{itemize}
\item The archimedean place $v = \infty$ (real numbers, $\mathbb{R}$):
  large-scale classical spacetime.
\item The non-archimedean places $v = p$ (p-adic numbers, $\Qp$):
  the arithmetic microstructure.
\end{itemize}

The adelic product formula $\prod_v |x|_v = 1$ constrains the
total vacuum energy. If certain primes are ``partially muted''
by cosmological evolution (e.g., by the dilution of
$p$-adic modes during inflation), the effective vacuum energy
would be
\begin{equation}
  \rho_{\mathrm{dark}} \propto \zeta(-3) \cdot
  \prod_{p \in S_{\mathrm{muted}}} (1 - p^{3}),
\end{equation}
which can take a very small value if the muted prime set
$S_{\mathrm{muted}}$ is chosen appropriately.

The observed dark energy density
$\rho_{\mathrm{dark}} \sim 10^{-9}\,\mathrm{J/m^3}$ constrains
the muted prime set. This is a quantitative prediction of the
arithmetic framework that could, in principle, be tested.

% ============================================================================
\section{Fundamental constants from arithmetic geometry}
\label{sec:constants}
% ============================================================================

If spacetime is $\Spec(\mathbb{Z})$, then dimensionless physical
constants should be computable from its arithmetic-geometric
invariants.

\begin{conjecture}[Arithmetic constants]
\label{conj:constants}
The dimensionless constants of physics (fine structure constant
$\alpha$, weak mixing angle $\theta_W$, Yukawa couplings, etc.)
are determined by the arithmetic geometry of $\Spec(\mathbb{Z})$:
zeta special values, K-group orders, regulator volumes, and
class numbers of number fields.
\end{conjecture}

Partial evidence from Connes' NCG Standard Model~\cite{CCM2007}:
the spectral action $\Tr(f(D/\Lambda))$ on $M^4 \times F$
determines the Standard Model coupling constants at unification
scale as functions of the parameters of $F$.
If $F$ has BC structure (Conjecture~\ref{conj:arith}), these
parameters are arithmetic invariants.

Specifically:
\begin{itemize}
\item $\zeta(-1) = -1/12$ appears in the 1-loop vacuum energy,
  related to the trace anomaly and conformal symmetry breaking.
\item $\zeta(-3) = 1/120$ determines the 3D Casimir energy,
  related to the cosmological constant.
\item $\zeta(2) = \pi^2/6$ and $\zeta(4) = \pi^4/90$ appear in
  black-body radiation (Stefan-Boltzmann law), connecting
  thermodynamics to $\zeta$ values.
\item $K_2(\mathbb{Z}) = \mathbb{Z}/2$ (Milnor) may encode
  the discrete symmetries (C, P, T) of the Standard Model.
\end{itemize}

% ============================================================================
\section{The multiverse as localization}
\label{sec:multiverse}
% ============================================================================

If our universe corresponds to $\Spec(\mathbb{Z})$, what
corresponds to other possible universes?

\begin{conjecture}[Arithmetic multiverse]
\label{conj:multiverse}
The ``multiverse'' is the set of localizations of $\mathbb{Z}$:
\begin{equation}
  \mathcal{M} = \{\Spec(\mathbb{Z}[S^{-1}]) : S \subseteq \mathrm{Primes}\}.
\end{equation}
Each choice of $S$ (a set of primes to ``mute'') defines a
universe with a different \'etale topos, different internal logic,
and therefore different physical laws.
\end{conjecture}

The extreme cases:
\begin{itemize}
\item $S = \emptyset$: $\Spec(\mathbb{Z})$. Our universe. All
  primes active. Full WEC. Standard physics.
\item $S = \{p\}$: $\Spec(\Zp)$. A universe missing one prime.
  WEC violated for $p$. Warp drive possible.
\item $S = \mathrm{all\ primes}$: $\Spec(\mathbb{Q})$. A universe
  with no primes. All Euler factors removed.
  $\zeta_{S}(s) = 1$ (trivial). Zero vacuum energy.
  No particles, no forces, no structure.
\end{itemize}

This provides a natural ordering of universes:
more primes muted $\Rightarrow$ simpler vacuum structure
$\Rightarrow$ fewer physical phenomena.
Our universe ($S = \emptyset$) is the ``maximal'' universe,
with all primes active and the richest possible physics.

\textbf{The anthropic principle, arithmetized.}
The question ``why does our universe have these physical laws?''
becomes ``why is $S = \emptyset$?'' --- or equivalently,
``why are all primes active?''
A possible answer: $S = \emptyset$ is the only choice that
supports sufficient complexity for observers (since muting primes
removes degrees of freedom needed for chemistry and biology).

% ============================================================================
\section{Observational tests}
\label{sec:tests}
% ============================================================================

\begin{table}[h]
\centering
\caption{Observational signatures of arithmetic spacetime.}
\label{tab:tests}
\small
\begin{tabular}{@{}p{3cm}p{4cm}p{4.5cm}@{}}
\toprule
Prediction & Observable & Current status \\
\midrule
$\Lambda = \zeta(-3) \times (\text{Planck})$ & Cosmological constant
  & Discrepant by $\sim 10^{120}$; may be resolved by partial muting \\
BC phase transition at Big Bang & CMB power spectrum with
  $\zeta$-function features & Not yet searched for \\
Dark energy $\propto \prod (1-p^3)$ & Equation of state $w(z)$
  evolution & Consistent with $w = -1$ (current data) \\
Constants from $\zeta$ values & Predictions for $\alpha, \theta_W$
  at unification & Requires explicit computation within NCG \\
Riemann zeros in CMB & Oscillations at positions of $\zeta$ zeros
  & Tantalizing claims exist~\cite{Wolf2010} but unconfirmed \\
\bottomrule
\end{tabular}
\end{table}

The most immediately testable prediction is the CMB signature
of the BC phase transition: oscillatory features in the
primordial power spectrum correlated with the Riemann zeros.
This can be searched for in existing Planck satellite data.

% ============================================================================
\section{Discussion}
\label{sec:discussion}
% ============================================================================

\subsection{Falsifiability}

Each conjecture in this paper is, in principle, falsifiable:
\begin{itemize}
\item Conjecture~\ref{conj:cc}: if the cosmological constant
  cannot be expressed in terms of $\zeta$ values and
  arithmetic invariants, the premise fails.
\item Conjecture~\ref{conj:bigbang}: if the CMB shows no
  trace of $\zeta$-function structure, the BC phase
  transition interpretation is disfavored.
\item Conjecture~\ref{conj:constants}: if the NCG spectral
  action fails to predict any measured constant, the
  arithmetic program loses predictive power.
\end{itemize}

\subsection{Relation to other approaches}

The arithmetic spacetime hypothesis intersects with several
established research programs:
\begin{itemize}
\item \textbf{$p$-adic cosmology} (Volovich 1987,
  Dragovich~\cite{Dragovich2017}): uses $p$-adic numbers in
  cosmological models. Our framework provides a unified
  (adelic) setting.
\item \textbf{Spectral geometry} (Connes~\cite{Connes1996}):
  derives physics from spectral data. We add the arithmetic
  ($\Spec(\mathbb{Z})$) perspective.
\item \textbf{Amplituhedron / motivic periods}
  (Arkani-Hamed~\cite{ArkaniHamed2014}, Brown~\cite{Brown2015}):
  scattering amplitudes involve periods of mixed motives.
  Our framework connects these to the vacuum structure
  via the motivic cohomology of $\Spec(\mathbb{Z})$.
\end{itemize}

% ============================================================================
\section{Conclusion}
\label{sec:conclusion}
% ============================================================================

If spacetime is $\Spec(\mathbb{Z})$, then:
\begin{enumerate}
\item The cosmological constant is $\zeta(-3)$ in natural units,
  not a cutoff artifact.
\item The Big Bang is a Bost-Connes phase transition where
  primes differentiate.
\item Dark energy is the Casimir energy of the arithmetic scheme.
\item Fundamental constants are arithmetic-geometric invariants.
\item The multiverse is the lattice of localizations of $\mathbb{Z}$.
\end{enumerate}

None of these are proven. All are falsifiable. Together, they
form a coherent cosmological narrative in which number theory
is not an abstract mathematical discipline but the science of
the universe's fundamental structure.

The experimental program of~\cite{paper_A} --- testing the
Euler product structure of vacuum energy in a superconducting
circuit --- is the first step toward confirming or refuting
this picture.

\begin{thebibliography}{20}
\bibitem{paper_A} Wright Brothers, companion paper A (2026).
\bibitem{paper_B} Wright Brothers, companion paper B (2026).
\bibitem{paper_C} Wright Brothers, companion paper C (2026).
\bibitem{Weinberg1989} S.~Weinberg, Rev. Mod. Phys. \textbf{61}, 1 (1989).
\bibitem{CCM2007} A.~H.~Chamseddine, A.~Connes, and M.~Marcolli, Adv. Theor. Math. Phys. \textbf{11}, 991 (2007).
\bibitem{Connes1996} A.~Connes, Commun. Math. Phys. \textbf{182}, 155 (1996).
\bibitem{Dragovich2017} B.~Dragovich et al., $p$-Adic Numbers Ultrametric Anal. Appl. \textbf{9}, 87 (2017).
\bibitem{ArkaniHamed2014} N.~Arkani-Hamed and J.~Trnka, JHEP \textbf{2014}, 030 (2014).
\bibitem{Brown2015} F.~Brown, Commun. Math. Phys. \textbf{338}, 1 (2015).
\bibitem{Wolf2010} M.~Wolf, arXiv:1002.1891 (2010).
\end{thebibliography}

\end{document}
